\section{Two inherited Gram forms and their common label dictionary}
\label{sec:tpc58-two-gram-dictionary}

The comparison problem starts with two exact identities from TPC--45,
not with two names for the same quadratic form.  We restate the identities
on one side of the equality column, retain the shaped terminal coefficient,
and insert the two successive quotient maps that are needed later.  The
right side is obtained by interchanging the two source systems; the full
column is their orthogonal direct sum.  No estimate in this section mixes
the two sides before that direct sum is taken.

\subsection{The very-high cutoff and the baseline atoms}

Fix the one physical intercept
\begin{equation}
 h_0\in\mathbb Z\setminus\{0\}.
 \label{eq:tpc58-fixed-intercept-dictionary}
\end{equation}
Let \(N_\pm,D_\pm,J_+\), and \(\Delta_M\) be the support endpoints in the
TPC--45 equality packet.  The very-high cutoff \(y=y_X\) is required to
satisfy both inherited conditions
\begin{align}
 y&>\max\{|h_0|,D_+,J_+,\Delta_M,\sqrt{N_+}\},
 \label{eq:tpc58-cutoff-geometric-condition}\\
 \frac{D_+N_+}{y}&<D_-N_-.
 \label{eq:tpc58-cutoff-label-separation}
\end{align}
At the endpoint one may take \(y=X^{267/400+\eta}\) for fixed
\(0<\eta<133/400\), but no exponent calculation below uses more than
\eqref{eq:tpc58-cutoff-geometric-condition}--
\eqref{eq:tpc58-cutoff-label-separation}.

This endpoint choice changes the admissible cofactor band.  On the physical
equality support,
\begin{equation}
 pr=mj+h_0\le N_+=X^{1+o(1)},
 \qquad p>y,
 \label{eq:tpc58-endpoint-product-ceiling}
\end{equation}
and therefore
\begin{equation}
 \boxed{r\le \frac{N_+}{y}
 =X^{133/400-\eta+o(1)}.}
 \label{eq:tpc58-compatible-cofactor-ceiling}
\end{equation}
For fixed \(\eta>0\), this is eventually disjoint from the inherited block
\(R\asymp J=X^{133/400+o(1)}\).  Thus every identity below is exact on the
declared very-high face, but applying it to endpoint energy requires a
separate lower transfer into a compatible smaller \(r\)-band.  The old
\(R\asymp J\) block may otherwise make the identities empty truths; no
occupancy statement is inferred from the cutoff alone.

Write \(i\in\cI\) for one native output coordinate.  On the left side,
\(i=(L,\gamma,j)\); on the right side, \(i=(R,\alpha,j)\).  TPC--45
associates to each actually occurring squarefree equality label \(u\) a
stripped coefficient \(b_{i,u}\in\C\).  The word \emph{stripped} means that
the M\"obius character \(\mu(u)\) is not contained in \(b_{i,u}\); the
 literal stripped source amplitude, mask, smooth profiles, and shaped
 terminal weight are retained.  We use the canonical active label set
\begin{equation}
 \cU_i^\circ:=\{u:b_{i,u}\text{ is not identically zero on the declared
 baseline carrier}\}.
 \label{eq:tpc58-baseline-native-labels}
\end{equation}
For one distinguished coefficient vector some entries on this fixed
carrier may vanish.  They remain zero coordinates; they do not change the
definition of the carrier or create a new representation.

The cutoff conditions give the unique factorization
\begin{equation}
 u=sq,
 \qquad \Pplus(s)\le y,
 \qquad q=1\ \text{or}\ q=p\in\cP_y,
 \label{eq:tpc58-low-high-factorization}
\end{equation}
where \(\cP_y\) is the set of occurring primes \(p>y\).  If \(q=p\), then
\(s=dr\) on the literal equality support.  For each fixed pair \((i,p)\)
there is at most one source row on the chosen side.  This last statement is
the inherited fixed-\((j,p)\) singleton incidence; it does not say that a
fixed \((i,s)\) contains at most one terminal prime.

\subsection{The native affine vectors}

In the native output Hilbert space
\begin{equation}
 \cH_{\rm nat}:=\ell^2(\cI)
 \label{eq:tpc58-native-space}
\end{equation}
define
\begin{align}
 (v_0)_i
 &:=\sum_{\substack{u\in\cU_i^\circ\\u_{>y}=1}}
       \mu(u)b_{i,u},
 \label{eq:tpc58-native-base-vector}\\
 (v_p)_i
 &:=\sum_{\substack{s:\,sp\in\cU_i^\circ}}
       \mu(s)b_{i,sp},
 \qquad p\in\cP_y.
 \label{eq:tpc58-native-high-vectors}
\end{align}
The sum in \eqref{eq:tpc58-native-high-vectors} has at most one nonzero
term for fixed \((i,p)\), but writing the sum keeps the output map visible.
At the physical all-minus prime-sign point the exact native column is
\begin{equation}
 \boxed{
 S_{\rm nat}(h_0)=v_0-\sum_{p\in\cP_y}v_p,}
 \qquad
 \cN_{\rm nat}(h_0)=
 \left\|v_0-\sum_{p\in\cP_y}v_p\right\|_{\cH_{\rm nat}}^2.
 \label{eq:tpc58-native-affine-column}
\end{equation}
The minus sign is the physical value of the unique prime \(p>y\).  It is
part of the affine identity and cannot be discarded when the base--high
cross term is later formed \citep{WangTPC45}.

\subsection{The residual-label Gram and the M\"obius gauge}

The high residual Hilbert space is
\begin{equation}
 \cH_{\rm res}^{H}
 :=\ell^2\{(i,s):i\in\cI,\ \Pplus(s)\le y,\
                   \text{ and }sp\in\cU_i^\circ
                   \text{ for some }p\in\cP_y\}.
 \label{eq:tpc58-high-residual-space}
\end{equation}
For \(p\in\cP_y\), the TPC--45 residual vector is the \emph{stripped}
vector
\begin{equation}
 (x_p)_{i,s}:=
 \begin{cases}
  b_{i,sp},&sp\in\cU_i^\circ,\\
  0,&sp\notin\cU_i^\circ.
 \end{cases}
 \label{eq:tpc58-stripped-residual-vector}
\end{equation}
Define the residual M\"obius gauge and the unsigned residual-label erasure
by
\begin{align}
 (\Uop_\mu z)_{i,s}&:=\mu(s)z_{i,s},
 \label{eq:tpc58-mobius-gauge}\\
 (\Eop z)_i&:=\sum_s z_{i,s}.
 \label{eq:tpc58-dictionary-unsigned-erasure}
\end{align}
Squarefreeness makes \(\Uop_\mu\) unitary.  The definitions give the
coordinatewise identity
\begin{equation}
 \boxed{\Eop\Uop_\mu x_p=v_p.}
 \label{eq:tpc58-primewise-erasure-identity}
\end{equation}
Consequently
\begin{equation}
 \|x_p\|_{\cH_{\rm res}^{H}}
 =\|\Uop_\mu x_p\|_{\cH_{\rm res}^{H}}
 =\|v_p\|_{\cH_{\rm nat}}.
 \label{eq:tpc58-primewise-diagonal-equality}
\end{equation}
The last equality uses the fixed-\((i,p)\) singleton incidence.  It is not
an assertion that \(\Eop\) is an isometry on a sum of different prime
columns.

To fix the inherited global sign once and for all, put
\begin{equation}
 \boxed{
 z_H:=-\sum_{p\in\cP_y}x_p,
 \qquad
 A_H:=\Uop_\mu z_H.}
 \label{eq:tpc58-high-residual-sign-convention}
\end{equation}
Thus \(z_H\) is the stripped resolved high vector, whereas \(A_H\) is its
physical residual-M\"obius version.  Equations
\eqref{eq:tpc58-primewise-erasure-identity} and
\eqref{eq:tpc58-high-residual-sign-convention} imply
\begin{equation}
 \boxed{
 \Eop A_H=-\sum_{p\in\cP_y}v_p.}
 \label{eq:tpc58-high-native-erasure}
\end{equation}
This is the sign convention used throughout TPC--58.

The low-sign-averaged identity of TPC--45 is
\begin{equation}
 \boxed{
 \mathbb E_{q\le y}Z_0(\varepsilon_{\le y},-\one_{>y})
 =\cD_0-\cD_H+\|A_H\|_{\cH_{\rm res}^{H}}^2.}
 \label{eq:tpc58-low-averaged-residual-gram}
\end{equation}
Here \(\cD_0\) is the raw stripped atomic diagonal and \(\cD_H\) is its
very-high-prime part.  The equality
\(\|A_H\|=\|z_H\|=\|\sum_px_p\|\) uses only the common diagonal gauge.
In particular, the last term in
\eqref{eq:tpc58-low-averaged-residual-gram} already contains the coherent
cross terms between different terminal primes that share the same
\((i,s)\).  It is not the terminal atomic diagonal
\(\sum_p\|x_p\|^2\).

\begin{proposition}[Exact two-Gram dictionary]
\label{prop:tpc58-two-gram-dictionary}
Under \eqref{eq:tpc58-cutoff-geometric-condition}--
\eqref{eq:tpc58-cutoff-label-separation}, the native and residual
very-high faces are related by
\begin{align}
 \sum_p\|v_p\|^2
 &=\sum_p\|x_p\|^2=\cD_H,
 \label{eq:tpc58-two-gram-common-diagonal}\\
 \left\|-\sum_pv_p\right\|^2
 &=\|\Eop A_H\|^2,
 \label{eq:tpc58-two-gram-native-high}\\
 \left\|\sum_px_p\right\|^2
 &=\|A_H\|^2.
 \label{eq:tpc58-two-gram-residual-high}
\end{align}
For \(p\ne q\), however,
\begin{align}
 \langle v_p,v_q\rangle_{\cH_{\rm nat}}
 &=\sum_i\sum_{s,t}
   \mu(s)\mu(t)b_{i,sp}\overline{b_{i,tq}},
 \label{eq:tpc58-native-cross-product}\\
 \langle x_p,x_q\rangle_{\cH_{\rm res}^{H}}
 &=\sum_i\sum_s b_{i,sp}\overline{b_{i,sq}}.
 \label{eq:tpc58-residual-cross-product}
\end{align}
Thus the residual Gram retains only \(s=t\), whereas the native Gram also
contains every \(s\ne t\) alias pair sharing the same output \(i\).
\end{proposition}

\begin{proof}
The common diagonal is
\eqref{eq:tpc58-primewise-diagonal-equality} summed over \(p\).
Equations \eqref{eq:tpc58-two-gram-native-high} and
\eqref{eq:tpc58-two-gram-residual-high} follow from
\eqref{eq:tpc58-high-residual-sign-convention}--
\eqref{eq:tpc58-high-native-erasure}.  Expanding the two inner products
gives \eqref{eq:tpc58-native-cross-product}--
\eqref{eq:tpc58-residual-cross-product}.
\end{proof}

The proposition is an exact L1 dictionary for the inherited face.  It does
not give either inequality between
\(\|\Eop A_H\|^2\) and \(\|A_H\|^2\).  Nor may the nonnegative base term
\(\cD_0-\cD_H\) in
\eqref{eq:tpc58-low-averaged-residual-gram} be replaced by
\(\|v_0\|^2\): the former is obtained after low-sign orthogonalization,
while the latter retains the native base--base cross terms.

\subsection{The two quotient maps are not interchangeable}

For a literal terminal atom write its prime-resolved key as
\begin{equation}
 (i,s,p),
 \qquad i=(L,\gamma,j)\ \text{or}\ (R,\alpha,j),
 \qquad s=dr.
 \label{eq:tpc58-prime-resolved-key}
\end{equation}
The two maps in the physical chain are
\begin{equation}
 \boxed{
 (i,s,p)\longmapsto(i,s)
 \longmapsto i.}
 \label{eq:tpc58-two-grouping-maps}
\end{equation}
The first arrow coherently sums terminal primes in a fixed residual label;
it is the \(p\to(i,s)\) grouping studied in TPC--55--57
\citep{WangTPC55,WangTPC56,WangTPC57}.  The second arrow is the unsigned
erasure \(\Eop\); it aliases different residual labels \(s\) in the same
native output coordinate.  A singleton for the first arrow need not be a
singleton for the second.  In particular, a lower square proved after the
first grouping is an input to the native-alias problem, not its solution.

All displayed sums use raw counting measure on their stated canonical
coordinates.  Prime-resolved injectivity is used before the first arrow,
and the \(p\)-sum is formed exactly once in each \((i,s)\).  Duplicating a
prime-resolved atom, adjoining an ambient affine-prime solution, or
inserting cancelling null pairs would change the Hilbert space and is not
part of this dictionary.
